% r03_h6w_b: proof of Theorem rot:thm:H6W, part 2: graphs and the census of the source subgraphs, the restricted form of King's renormalisation theorem, and the Cauchy estimate at the frozen wall modulo the field-free constant.

\subsubsection*{Graphs, degrees and the census of the source subgraphs}

King's perturbative part $P^{(r)}_k(tG,h)$ is a finite sum of connected renormalised graphs $H$ \cite[(3.29), (3.37)]{K1} with the vertices \cite[(3.18)--(3.26)]{K1}; the number of graphs depends on $\bar n$ only, each graph is multilinear in its source insertions, and its value on $\supp\chi_r$ is $\sum_{\{y\},\{z\}}E^{(r)}(H;\{y\},\{z\})\times(\text{external fields at }\{y\})$, with a kernel that depends on the coarse background $A^{(r)}$ through its lines $G_r(A^{(r)})$ \cite[(3.17)]{K1}, so that a graph without external legs still depends on the background. Degrees are King's \cite[(3.30)--(3.31)]{K1}, as in \Cref{rot:lem:S1}, computed on the point-split graph (the two summation points of a block product as vertices, as in \Cref{rot:cor:OF}), with two conventions. The curvature vertex \cite[(3.20)]{K1} is counted with $D_v=\frac52$, the lattice curl of its photon leg having been summed by parts onto the smooth $G$ (exact on the torus, $|d^*_\eta G|\le C\|\nabla G\|_\infty$); with the curl charged to the vertex, $D_v=\frac32$ and the tree with two curvature vertices would have degree $0$, so the summation by parts is necessary. The $h$-vertex \cite[(3.21)]{K1} has $D_v=2$ (\Cref{rot:lem:S1}). Every connected subgraph containing an $h$-vertex and a curvature vertex therefore has $D\ge2+\frac52-3=\frac32>0$, and the smallest vacuum graph with two curvature vertices and two current vertices has $D=5$.

Each graph value is split at zero coarse fields, $E(H)=E(H)|_0+[E(H)-E(H)|_0]$, where $|_0$ means $(A_r,\phi_r)=(0,0)$, and we put $E_N(H)\defeq E(H)$ if $H$ has an external scalar leg or a curvature vertex and $E_N(H)\defeq E(H)-E(H)|_0$ otherwise: the field- or source-dependent part. The field-free, $G$-free values $E(H)|_0$ of the remaining graphs ($h$-linear or $h$-nonlinear vacuum graphs without curvature vertex, and the $h$-free vacuum energy) are numbers depending on $k$, $r$, $h$ only; nothing is claimed about their convergence, and for the order-$e^2$ $h$-linear family it fails (\Cref{rot:rem:log}: the Wick constant drifts by $c_*\ln5$ per step).

\begin{lemma}[the census of the source subgraphs]\label{rot:lem:H6census}
Expand every graph in the coarse background as in \cite[(3.43)--(3.55)]{K1} (\Cref{rot:lem:CB}(1) for one flow): $E(H)$ becomes a finite sum of graphs $\tilde H$ with lines at zero background, explicit background legs $\tilde A$ at block points, and remainder factors, and $E(H)|_0$ is the core $\tilde H_0$ without background legs. Then:
\begin{enumerate}[label=\textup{(\alph*)},leftmargin=*]
\item Every $\tilde H$ entering $E_N(H)$ is connected and has an external leg (scalar, background or $dG$) or a curvature vertex.
\item For such $\tilde H$, every connected ordered subgraph $\gamma$ (lines with scale indices below $r$) has strictly positive King degree or vanishes identically. In detail: if $\gamma$ contains a block point, $D(\gamma)\ge1$ ($\ge2$ with a source vertex); if $\gamma$ contains an $h$-vertex, $D(\gamma)\ge\frac12$ unless $\gamma$ is a vacuum graph with exactly one $h$-vertex and no curvature vertex (then $\gamma=\tilde H$, excluded by (a)) or $\gamma$ is the $h$-loop through a single current vertex \cite[(3.22)]{K1} with its photon leg external ($D=0$: the current tadpole of \Cref{rot:lem:S1} with the $h$-insertion), which vanishes identically at zero background ($\Tr q=0$); if $\gamma$ contains a curvature vertex and no $h$-vertex, $D(\gamma)\ge\frac12$ unless $\gamma$ is the curvature vertex joined by a photon line to a current tadpole ($D=0$), which vanishes identically at zero background for the same reason; if $\gamma$ is source-free, $D(\gamma)>0$ after Ba{\l}aban's resummed rewriting \cite[(3.5)]{Ba3}, the input of \Cref{rot:cor:OF}.
\item The rewriting of the source-free subdivergences inside a graph that carries source vertices (Ba{\l}aban's Taylor subtraction at the subdivergence with its counterterm partner, \cite[(3.5)]{Ba3}, \cite[(3.83)--(3.85)]{K1}) does not lower the degree of a prefix component that contains a source vertex below the value counted in (b).
\item The only connected subgraphs containing a source vertex with King degree $\le0$ in the raw count are the field-free $h$-linear vacuum graphs of order $\le2$ in the couplings, the family $\cT_h$ of \Cref{rot:lem:S1} (which have no external leg, hence occur only as whole graphs), and the two vanishing exceptions of (b).
\end{enumerate}
Consequently the minimum $\alpha_0$ of the degrees of the non-vanishing ordered subgraphs over the finitely many graphs and orderings is strictly positive.
\end{lemma}
\begin{proof}
(a) The core $\tilde H_0$ is the same in $E(H)$ and $E(H)|_0$ and cancels in $E_N(H)$ when $H$ has no external scalar leg and no curvature vertex; every other $\tilde H$ carries a background leg, an external $dG$ line or a remainder factor (the remainders are bounded separately, \Cref{rot:lem:CB}(2)). (b) The degree arithmetic is that of \Cref{rot:lem:S1}: inserting an $h$-vertex into a scalar line adds $2$; source-free connected subgraphs with an external leg have $D\ge-1$ except the current tadpole ($D=-2$), so a subgraph with an $h$-vertex and an external leg has $D\ge1$ except the $h$-loop through one current vertex ($D=0$); a vacuum subgraph with exactly one $h$-vertex has $D\le0$ only at order $\le2$ in the couplings, the exhaustive enumeration of \Cref{rot:lem:S1}, and it is then the whole graph; with two or more $h$-vertices $D\ge1$. For a curvature vertex, which attaches by a photon leg and contributes $\frac52$, the same enumeration leaves the one exception named. Both exceptions vanish at zero background because the current one-point function of the loop vanishes, $\Tr q=0$, exactly as for the current tadpole in \Cref{rot:lem:S1}. A component containing a block point has degree $\ge1$ by the point-split count of \Cref{rot:cor:OF} (every component containing a block point has Ba{\l}aban's degree $\ge1$, \cite[(2.1)--(2.2)]{Ba3}), and a source vertex adds at least $1$ more. (c) Ba{\l}aban's Taylor subtraction at a source-free subdivergence $S$ raises the degree of $S$ by $\theta_S$ and lowers by $\theta_S$ the degree of the external line of $S$ that receives the derivative; the construction, run on the whole graph with the source vertices counted as vertices, rewrites $S$ only in orderings in which every line of $S$ is finer than every external line of $S$. Let $\gamma$ be a component of an ordered subgraph. If the line that receives the derivative lies in $\gamma$, then every line of $S$, being finer, lies in the same ordered subgraph, and $S$, connected to that line, lies in $\gamma$: the two changes cancel. Otherwise the degree of $\gamma$ can only increase. By induction over the forest of subdivergences, finest first, no component falls below its value in (b); the source vertices, which are never inside a rewritten $S$, keep their contributions $2$ and $\frac52$. (d) is read off the same arithmetic.
\end{proof}

\subsubsection*{The restricted renormalisation theorem}

King's renormalisation theorem \cite[Thm.~3.5]{K1} asserts the per-site bound (3.38) and the convergence (3.39) of the graph kernels for all graphs. As printed it is false on the order-$e^2$ $h$-linear field-free family (\Cref{rot:rem:log}), and the convergence of the field-free, $G$-free constants is not needed here. What is used is the following restricted form.

\begin{proposition}[the restricted renormalisation theorem]\label{rot:prop:H6R}
Let $h$ be in the family of \Cref{rot:lem:H6conv}, $t$ real, and let $\gamma_*,\beta_*,\sigma$ be the exponents of \cite[Thm.~3.4]{K1}.
\begin{enumerate}[label=\textup{(\roman*)},leftmargin=*]
\item \textup{(convergence)} For every graph $H$ and every $k$, with $r=r(k)$, the field- or source-dependent parts obey, after the sum over the positions, for any $\gamma'<\gamma_*$, $\beta'>\frac72$, $\sigma'<\sigma$,
\[
\sum_{\text{positions}}\big|E^{(r+1)}_{N,k+1}(H)-E^{(r)}_{N,k}(H)\big|\times|\text{external fields}|\le C\big[5^{-\gamma'r}\delta_r^{-\beta'}+\delta_r^{\sigma'}\big]|T|\quad\text{on }\supp\chi_r,
\]
the two kernels living on the common unit lattice $T^{(r)}$ of $\Xk$ at depth $r$ and $X_{k+1}$ at depth $r+1$; $C=C(t,G,h_0,|T|)$ is independent of $k$ and locally uniform in $t$.
\item \textup{(per-site bound)} For every graph $H$ with at least one $h$-vertex, other than the degree-$0$ members of the order-$e^2$ $h$-linear field-free family,
\[
|E^{(r)}_k(H;\{y\},\{z\})|\le C(1+k-r)^{p_0}\,\delta_r^{\varrho_H}\,p(\delta_r)^{n_1}\,e^{-\delta\dist(\{y\},\{z\})},\qquad\varrho_H\ge0,
\]
and for the degree-$0$ members of that family the same bound with the extra factor $(1+r)$, which is necessary there. For a graph at non-zero background the bound is read through the expansion of \Cref{rot:lem:H6census}: the core at zero background falls under the second clause if $H$ is in that family and under the first otherwise, and every other term has an external leg and falls under the first. $C$ is linear in $\sup|h|$ for a distinguished $h$-vertex and depends on the other sources through sup norms; $p_0$, $n_1$ depend on $\bar n$ only.
\end{enumerate}
\end{proposition}
\begin{proof}
King's proof of (3.39) \cite[pp.~660--662]{K1} reduces it to the convergence of a single graph with one line replaced by the difference of the two flows' lines (3.56), and that to \cite[Prop.~3.6]{K1}, the one-flow estimate for a renormalised family with one element weakened by a small degree, which King asserts on p.~665 ``by redoing the analysis of Sect.~3.4'' and which is \Cref{rot:cor:OF}. \Cref{rot:cor:OF} applies to a family exactly when every connected ordered subgraph of its non-vanishing graphs has strictly positive degree after Ba{\l}aban's rewriting: a zero-degree ultraviolet subgraph containing the replaced line would give $\sum_j(5^j\eta)^{-\gamma}\sim5^{\gamma r}$ and cancel the gain $5^{-\gamma r}$. By \Cref{rot:lem:H6census} this holds, with the minimum $\alpha_0>0$, for every $\tilde H$ entering $E_N(H)$. The replaced line is a fine line of \cite[(3.17)]{K1} with scale index below $r$; the lines with indices between $r$ and $k$ occur only on the counterterm lines of the $\dm$ graphs, never at a source vertex, so the counterterm-attachment clause \cite[(3.78)]{K1} and the one-vertex attachment are the source-free statements and are unchanged by the presence of the sources. The remaining steps are King's: the kernel bounds \cite[Thm.~3.3, Props.~3.7--3.9]{K1} for the lines and the fine-field reconstructions, printed for the torus convention at block factor $5$ (\Cref{thm:a06:BALA}); the sampling of the sources $x'\to x$ (p.~665), which costs $C\eps_k\|\nabla h\|_\infty+C\eps_k\|\nabla G\|_\infty\le C5^{-\gamma r}$; the tree-level $h$-term $\delta_r^2\sum\eta^3h|\phi^{(r)}|^2$, which has no internal line and is compared directly through the convergence of the fine-field reconstructions $a_rG_rQ_r^*$ \cite[Thm.~3.3, \S3.6]{K1}; the non-polynomial terms of (3.14) by \cite[\S3.6, (3.89)--(3.98), Prop.~3.10]{K1}; the Taylor remainders beyond order $\bar n$, $C\delta_r^\sigma|T|$ per flow (p.~665); and the position sum (3.40)--(3.42), which gives $|T|$ with one free position, the others tied by the exponential tree decay. The factor $(1+k-r)^{p_0}$ of \Cref{rot:cor:OF} is a power of $\log\delta_r^{-1}$ and is absorbed by $\gamma'<\gamma_*$, $\sigma'<\sigma$, and the field norms $(\lambda\delta_r)^{-s/4}p(\delta_r)^{n}$ of the external legs by $\beta'>\frac72$. This proves (i).

(ii) For a graph outside the degree-$0$ members of the order-$e^2$ $h$-linear field-free family, \Cref{rot:lem:H6census} gives strictly positive ordered degrees for every non-vanishing derived graph (the excluded vacuum case of the census is now the family itself), and \Cref{rot:cor:OF} without weakening gives the per-site bound with the factor $(1+k-r)^{p_0}$, King's ``possible addition of some power of $\ln(L^k\eps)^{-1}$'' (p.~666); the sup-norm dependence on the sources is King's (p.~664, 677), the distinguished vertex entering linearly. The two $h$-linear field-free members of order $0$ and $\lambda$ are bounded directly: the subtracted tadpole at zero coarse field is the infrared remainder $-N\sum\eta^3\delta_r^2h\,F_{r}(x,x)$ of the coincident propagator, bounded per site at unit scale, and the figure-eight minus its $\dm_1$ partner leaves a bubble of positive ordered degree times an infrared tadpole line, as in step 3 of the proof of \Cref{rot:cor:OF}. For the degree-$0$ members of the order-$e^2$ family (the exchange and seagull tadpoles with their $\dm_2$ partner) every proper ordered subgraph has $D>0$ (a proper component with the $h$-vertex is not a vacuum graph, the $h$-loop vanishes, and the source-free components are the self-energy and the photon tadpole, renormalised by the $\dm_2$ partner) and only the whole graph has $D=0$: the Taylor identity for the self-energy with its $\dm_2$ partner, Ba{\l}aban's general step at the whole graph with no final renormalisation (there is none to perform), as in King's worked case \cite[(3.81)--(3.87)]{K1}, makes every proper prefix component of degree $\ge\alpha$ and leaves the whole graph at degree $0$, so in the multiscale sum every geometric series keeps its constant except the last, $\sum_{j\le r-1}(5^j\eta)^0=r$: this is the factor $(1+r)$, the zero-degree logarithm of \Cref{rot:cor:OF}. It is necessary: the value of the family at a constant source is $u$ times the drift of \Cref{rot:rem:log}, of size $r$ at fixed $\delta_r$. The infrared counterterm lines are handled as in (i).
\end{proof}

\subsubsection*{The Cauchy estimate at the frozen wall}

\begin{proposition}[King's Cauchy estimate modulo the field-free constant]\label{rot:prop:H6I3}
Let $h=\hinf$ and $t\in\R$, and define the field-free constant $a^{(r)}_k(h)\defeq S^{(r)}_k(0,h)\big|_{A_r=\phi_r=0}$ (the whole value of the effective action at zero fields and zero curvature source) and $d_k\defeq a^{(r)}_k(h)-a^{(r+1)}_{k+1}(h)$, a number independent of $t$ and of $G$. Then on $\supp\chi_r$, with $\gamma',\beta',\sigma'$ as in \Cref{rot:prop:H6R}(i),
\[
\big|S^{(r)}_k(tG,h)-S^{(r+1)}_{k+1}(tG,h)-d_k\big|\le\eta_k|T|,\qquad\eta_k\defeq C\big[5^{-\gamma'r}\delta_r^{-\beta'}+\delta_r^{\sigma'}\big],
\]
$C$ independent of $k$ and locally uniform in $t$.
\end{proposition}
\begin{proof}
Both effective actions live on $T^{(r)}$ with the same $\chi_r$. Their Gaussian parts converge by \cite[Thm.~3.3]{K1} (the step King takes for (3.9) before the perturbative part), and their difference is a quadratic form in the fields plus a constant, the constant being part of $d_k$. In the perturbative parts, split every graph as $E(H)=E(H)|_0+E_N(H)$ when $H$ has no external scalar leg and no curvature vertex: the sum of the values $E(H)|_0$ over these graphs, at $G=0$, is by definition the field-free constant $a^{(r)}_k(h)$ minus the Gaussian constant, and the graphs with a curvature vertex and no external leg are included in $E_N$ and converge by \Cref{rot:prop:H6R}(i). So $S^{(r)}_k(tG,h)-S^{(r+1)}_{k+1}(tG,h)-d_k$ is a sum of differences $E^{(r+1)}_N(H)-E^{(r)}_N(H)$ over the finitely many graphs, each bounded by \Cref{rot:prop:H6R}(i), plus the two Taylor remainders, $\le C\delta_r^\sigma|T|$ each, plus the non-polynomial terms, included in \Cref{rot:prop:H6R}(i). This is King's \cite[(3.9)]{K1} without the field-free $h$-sector constant, which is exactly what is false in the printed form at $\int h\ne0$ (\Cref{rot:rem:log}: the order-$e^2$ $h$-linear constant drifts by $c_*\ln5\int h$ per step, so $d_k\not\to0$); $d_k$ is $t$- and $G$-independent by construction and cancels between numerator and denominator of every ratio below. Only one step, $n=1$, is used.
\end{proof}
